\chapter{Introduction}
\label{chap:intro}
% the word introduction should be a H1 or H2?
% the "\section{Background}" is autotagged as H2
Tonal languages are estimated to constitute 60 to 70 percent of the world's languages \citep{yip2002tone}, making tonal phonology a critical and unique subfield within phonology. One key characteristic distinguishing tonal phonology from segmental phonology is its autonomy: tones can operate independently of the segments with which they are associated. This autonomy has been captured by Autosegmental Phonology which uses a multi-tiered structure for tonal representation \citep{goldsmith1976autosegmental} and gives more accurate predictions and explanations than linear phonology \citep{clements1984autosegmental,oddenbook}. Computationally, processes in segmental phonology are characterized as sub-sequential and can be handled by a deterministic transducer \citep{heinz2013vowel} while some processes in tonal phonology such as unbounded tone spreading are found to be more complex and non-deterministic \citep{jardine16dissertation,jardineLocalNatureToneassociation2017}. However, beyond the formal properties of tones \citep{bird1990phonological,coleman1991no,kornai2018formal,jardine2015concatenation}, limited work has been done on tonal pattern learning, and more importantly, little effort has been made to capture the multi-linear properties of tones in the computational linguistics field \citep{albro1994amar,kornai2018formal}, despite their prevalence in the theoretical world.

In the proposed dissertation, I aim to provide a computational expression using formal logic to make tonal patterns and processes applicable for computational modeling and learning. While the previous attempts in tonal learning use a string representation \citep{chandleeComputationalAccountTone2019,oakden2020computational}, this work plans to develop a multi-linear structure faithful to linguistic representation that captures the autonomy of tones. Secondly, autosegmental phonology has achieved a great number of success in African tonal languages \citep{clements2010autosegmental,hyman2009representation}, but far less research has explored Asian tonal languages which are known for having intricate tone-segment interactions. Incorporating these interacting features into the modeling is another goal of this project. Finally, it is noteworthy that one possible reason that tonal phonology has not obtained the same level of discussion in computational linguistics is due to limited tools and resources. No algorithm to date has been able to effectively learn the input-output mapping of autosegmental representations, which are intrinsically graphs. Besides, tonal languages are often transcribed without tonal specifications (as in African language dictionaries) or are transcribed in their underlying representations (as in Mandarin Chinese). For those languages that do have sufficient data, there are usually varying representations or models \citep{bao1999structure,yip1989contour,sy1967phonological}. The proposed work aims to clarify these issues and hopes to develop more tools for tonal linguists to use in learning tonal patterns and processes. These goals will help to gain more insights into tonal phonology and its learnability, provide implications for theoretical generalizations, and get a better understanding of underrepresented languages with limited data.




\section{Thesis Structure}
	Chapter 3, ``Computation of Tonotactic Patterns,” is ongoing and is planned to be completed by the end of this year. Building on my second qualifying paper, this work involves expanding the existing algorithm to encode syllable weight and word boundaries in Hausa, and then apply it to Japanese Nouns. Corpora for both languages have already been gathered and are ready to be implemented into the modified BUFIA pipeline. My goal is to complete both experiments by the end of this semester and write up the chapter.

Chapter 4, ``Interaction between Tone and Other Features'' will be carried out roughly from the winter spanning over the spring semester of 2025. The winter time will be used for research, reading, and writing for the backgrounds. From February to April, I will start experimenting with perhaps two case studies from Shanghai Chinese, Mandarin, or Fuzhou. The plan for Chapter 3 depends on the progress of the two experiments. 


Chapter 5, “Tonal Process Learning,” will take the longest to complete, as there is limited research on the computation and modeling of this topic. I will need additional time to learn new material and conduct further reading. In the timeline, I have allocated May to September for this chapter, but I have also reserved October and November for both writing and addressing any unresolved issues that may arise during the completion of this chapter. 

The introduction chapter provides an overview, including three main points. First, tonotactic well-formedness and tonal processes differ from those at the segmental level, making it reasonable to adopt multi-tiered structures rather than string representations. Second, from a learning perspective, computational tonal phonology is still underexplored. Most existing studies use string representations; many focus on African languages, likely due to the convenience of obtaining tonal data; and very few learning models have been developed for either tonal processes or tonotactic patterns. These research gaps raise important questions regarding representation, patterns, and learning. These questions will be the main focus of this dissertation and will be addressed in the introduction chapter. Chapter One will also include logistical details, such as an outline of the dissertation.

Chapter 2 will serve as a Literature Review chapter and it will mainly cover two broader aspects around tonal phonology. First, as in the overview in this proposal, I will review the issues about how representation, grammar, and typology of tonal languages. Second, I will discuss how previous work regulates the formal properties of tonal representations using model theory or other mathematical methods. 

As the first major part of the dissertation, this chapter will investigate the tonotactic grammar of tonal languages by examining forbidden structures in all surface-true forms. In this chapter, I choose Hausa and Japanese as case studies to represent two typologically different tonal languages. Hausa uses a level tone system where H and L contrast and Japanese employs a pitch accent system. These patterns are argued to be local and their grammar can be learned by searching for the negative literal conjunctions. My goal is to represent the tonal patterns in model theory, apply a bottom-up learning algorithm to learn the grammar, and lastly compare the constraints found by the algorithm with previously reported generalizations, and see to what degree they match with each other.  

The frequently discussed tonal patterns in Hausa are generalized as No Rising Contour, Avoid LL sequence, No initial Contours, and No three-tone contours, but it is unclear to what extent these rules can explain the data, and there are many counterexamples against these generalizations.  In my QP2, I used a small Hausa corpus, from which I translated the orthographic representations into ARs and implemented BUFIA  on it. The experiment revealed 26 distinct ARs among 664 surface-true Hausa words and identified 7 AR structures (syllable number \(\leq 3 \), tone number \(\leq 3\)) missing from the data, which constitute the constraints in the tonal grammar. The computationally discovered constraints either match the ones previously proposed or provide more specific generalizations that better account for the data. 


\section{Motivations}

Effective tone modeling requires choosing appropriate representations and grammar. The motivation behind this work stems from the most common questions in tonal phonology, which address tonal representation, grammar, and learning.

\subsection{Representation}
The first question is what phonological representation should be used for tones. Under the influence of the classic framework of \textit{Sound Patterns of English} \citep{chomsky1968sound}, where phonological strings were divided into discrete feature matrices, tones were likewise treated as intrinsic properties of segments and represented with features indicating pitch ($\pm$ H) and contour shape ($\pm$ contour). Tonal changes are the result of changing feature values and insertion/deletion of segments. Contour comes into form due to two opposite pitch features are adjacent to each other. This approach, however, was soon argued to be inadequate for predicting correct attested patterns, and it failed to capture the commonly observed patterns among tonal languages \citep{hyman2008issues,oddenbook}. Therefore, ``remarkably little insight was produced by the SPE framework into the behavior of tones" \citep[p.8]{clements1984autosegmental}.

Abundant evidence from African tonal languages shows that when considering tonal forms, tonal units and their associations with timing units also need to be taken into account. For example, languages with ``restricted melodies," such as Mende \citep{leben2017tone}, have a fixed set of five melodies to any form regardless of the number of syllables of the form. These melodies ``must be abstracted away from the TBUs on which they are realized" \citep{voorhoeve1973safwa,hyman1974universals}. To capture this, Autosegmental Theory \citep{goldsmith1976autosegmental,clements1980vowel} enriches the geometry of phonological representation by proposing a multilinear structure: tones are represented on a separate autonomous tier and are related to segments on the adjacent tier through association lines governed by specific well-formedness conditions (WFC) (see also \citet{clements2010autosegmental}). Common tonal patterns such as contours or tone plateaus can be explained as multiple tones connecting to the same segment or one tone spanning across multiple segments. The insertion or deletion of one tier does not necessarily trigger the same operation on the other tier. The temporal rearrangements between tones and TBUs are handled by changing the association lines between two tiers. Autosegmental Theory has proven especially useful in explaining tonotactic patterns and tonal processes where linear phonological models often struggle \citep{hyman2011tone,oddenbook}. 

Feature matrices and Autosegmental Representations (ARs) are not the only approaches to tonal analysis. \citet{inkelas2013abc+} introduced \textit{Q-theory}, initially designed to account for complex diphthongs, which later expanded to contour tones. Under this framework, each tone (or segment) can be divided into three temporally ordered, quantized subsegments. This is motivated by the observation that most complex tones do not exceed three tonal levels and very rarely have four-tone contours \citep{shih2013subsegmental}. However, this approach has been argued to be logically equivalent to ARs in model theory and logical interpretations \citep{jardine2021qtheory}. \citet{kornai2018formal} also compared different ways of encoding string-based representations of ARs. Similarly, the two tonal models proposed by \citet{yip1989contour} and \citet{bao1999structure} are also shown to be notationally equivalent using formal logic \citep{oakdenNotationalEquivalenceTonal2020}.

Given that varying models and representations are argued to be logically equivalent, or at least translatable, formal logic is possible to represent their internal structure and facilitate learning explicitly. This is the motivation for this study. 

\subsection{Grammar}

Besides representations, the second question in tonal phonology concerns the type of grammar that should be used. The first distinction is whether the grammar is rule-based or constraint-based. Take the classic tonal mapping problem for instance. In a rule-based approach, tonal mapping requires a globally evaluated direction (leftward vs. rightward), along with other parameters such as morphological category and tone quality \citep{goldsmith1976autosegmental}, but this cannot account for those two-edge effects, and sometimes leads to wrong predictions. \citet{zoll2003optimal} argues that the directionality needs to be reconsidered, and proposes the optimal tonal mapping where directionality can be handled by the interaction between universal violable constraints: a morphological direction constraint \textsc{Align} and two quality-sensitive markedness constraints \textsc{Clash} and \textsc{Lapse}.

\citet{jardineLocalNatureToneassociation2017} argues that the grammar of tonotactic well-formedness consists of language-specific, inviolable local constraints. These constraints can be expressed locally over ARs, which are intrinsically graphs. A well-formed AR structure can be viewed as a graph that does not contain any forbidden subgraphs. Using Model Theory, each AR graph can be abstracted into a mathematical model with a unique signature to indicate its distinct internal structure. This method captures the locality and nonlinearity makes it possible to learn tonotactic patterns over ARs, and also explains the North Karanga Shona ``edge-in" pattern that Optimal Tonal Theory fails to generalize.
More importantly, this approach of searching subfactors to get grammar can be learned through The Bottom Up Factor Inference  Algorithm \citep[BUFIA,][]{chandleeLearningPartiallyOrdered2019}. 
\section{Contributions}

fukc you fuck you 